\section{Passive Meal-State Modeling}
\label{sec:meal-state-modeling}

\subsection{Motivation}
\label{subsec:meal-motivation}
AI-READI contains CGM, wearable signals, and static clinical descriptors, but not prospectively logged meals or carbohydrate doses \citep{aireadi2024}. The existing \texttt{predmeal\_flag} was therefore only a proxy scenario channel. We rebuilt the meal pathway as a weak-supervision problem and then asked a narrower question, whether any deployable meal feature improves held-out forecasting on the endpoints where a meal actually acts. The pipeline is summarized in Fig.~\ref{fig:meal-pipeline} and Table~\ref{tab:meal-pipeline-summary}.

\begin{figure}[t]
    \centering
    \includegraphics[width=\textwidth]{figures/meal_flags/meal_pipeline.pdf}
    \caption{Meal-flag improvement pipeline. The CGMacros teacher and the retrospective state decoder are diagnostic weak-supervision sources, the online Track A decoder and the causal student probability are the deployable candidates.}
    \label{fig:meal-pipeline}
\end{figure}

\begin{table}[t]
    \centering
    \caption{Meal-transfer pipeline summary.}
    \label{tab:meal-pipeline-summary}
    \begin{tabular}{lr}
\toprule
Quantity & Value \\
\midrule
Participants & 1,591 \\
Rows & 4,015,355 \\
Teacher probability finite fraction & 0.993 \\
Corrected teacher flag rate & 0.114 \\
Surviving old binary artifact rate & 0.042 \\
Pseudo-label positives & 507,830 \\
Pseudo-label negatives & 1,085,414 \\
Pseudo-label uncertain & 2,422,111 \\
Causal student validation AP & 0.863 \\
Causal student validation AUC & 0.928 \\
Retrospective decoded events & 53,024 \\
Retrospective decoded flag rate & 0.290 \\
Online prefix invariance max absolute difference & 0.000 \\
Online response-size validation macro-F1 & 0.767 \\
Independent evaluation events & 77,835 \\
\bottomrule
\end{tabular}

\end{table}

\subsection{Transfer of the CGMacros Meal Detector}
\label{subsec:meal-transfer-teacher}
The source audit recovered the CGMacros detector as a dilated CNN followed by a two-layer bidirectional LSTM and a per-time-step logit head over 72-sample, six hour windows of 5-minute CGM. Because the LSTM is bidirectional, the teacher probability at an anchor depends on CGM samples after the anchor, up to roughly six hours ahead. That window overlaps the forecast horizon, so the retrospective teacher reads future glucose about the exact quantity being forecast. We therefore treat the bidirectional teacher as a leakage diagnostic, not as a reachable target.

\subsection{Correction of the AI-READI Meal-Flag Reconstruction}
\label{subsec:meal-reconstruction}
The original downstream artifact preserved only a thresholded binary flag and discarded the teacher's continuous probability. Reconstructing overlap-averaged probabilities recovered finite values for 99.3\% of rows and a corrected teacher flag rate of 0.114, compared with 0.042 for the surviving old binary artifact. The reconstruction is a faithful recovery of the source signal, but the source signal is itself bidirectional, so neither the corrected probability nor the surviving \texttt{predmeal\_flag} is causal. Figure~\ref{fig:meal-teacher-reconstruction} shows the information loss of the old artifact.

\begin{figure}[t]
    \centering
    \includegraphics[width=0.82\textwidth]{figures/meal_flags/teacher_reconstruction.pdf}
    \caption{Corrected teacher reconstruction. The continuous teacher probability is available for nearly all rows, while the old binary artifact is merge-diluted.}
    \label{fig:meal-teacher-reconstruction}
\end{figure}

\subsection{Weak Target-Domain Supervision}
\label{subsec:meal-weak-supervision}
Weak labels were built from train-split quantiles, not hard-coded glucose thresholds. Two changes follow the leakage and pooling critique. First, insulin users are quarantined from meal-label training, because a decoded meal-like response in an insulin user often reflects unobserved insulin dynamics rather than a meal; 305,843 insulin rows are removed from the training labels and reported separately as low-confidence meal-like events. Second, the previously discarded uncertain rows are rescued with a soft, ordinal onset-distance target rather than dropped, because the pre-rise transition is the only region where a meal flag can lead the CGM trend; 389,623 of 2,422,111 uncertain rows now carry a nonzero soft target and stay in training. Response size remains a realized glucose-excursion proxy, not carbohydrate grams.

\subsection{Causal Student Model}
\label{subsec:meal-causal-student}
The target-domain student uses past and current glucose, wearable, clock, and static features only, and is now fit as a regressor on the soft onset-distance target over the quarantined, rescued training set. Teacher outputs and future glucose are used only during label construction. A good classification score does not imply a useful forecasting covariate, so the student is judged by the forecasting gate below, not by its pseudo-label AUC.

\subsection{Retrospective and Online Meal-State Decoding}
\label{subsec:meal-decoders}
A duration-constrained retrospective decoder converts the student probability into onset, rising, peak, and recovery states, and a strictly forward-only filter produces the deployable online states. The retrospective response-size proxy is physiologically ordered, mean peak rise increasing from 2.1 to 15.9 to 50.8 mg/dL across small, medium, and large events. The retrospective decoder is an event-discovery diagnostic because it uses full-segment summaries, the online decoder is the deployable candidate. Examples appear in Fig.~\ref{fig:meal-state-examples}.

\begin{figure}[t]
    \centering
    \includegraphics[width=\textwidth]{figures/meal_flags/state_examples.pdf}
    \caption{Strictly online decoded meal-state examples. State changes are produced by a forward recursion and do not use future glucose, future wearables, or teacher outputs at inference.}
    \label{fig:meal-state-examples}
\end{figure}

\subsection{Experimental Feature Representations}
\label{subsec:meal-feature-representations}

All setups share a common residual-correction framework. The \texttt{q50} median from the
pretrained SSMCGM-Stream model is used as a base forecast, and each candidate meal
representation $M_t^{(m)}$ is evaluated as an additive residual correction:
\begin{equation}
    \widehat{y}_{t+h}^{(m)}
    = \widehat{y}_{t+h}^{q50}
    + r_m\!\left(X_t^{\mathrm{base}},\,M_t^{(m)},\,h\right),
    \label{eq:meal-residual-framework}
\end{equation}
where the baseline feature vector is
\begin{equation}
    X_t^{\mathrm{base}}
    = \left[G_t,\;\Delta G_{15,t},\;\Delta G_{30,t},\;\Delta G_{60,t},\;
      \sin\theta_t,\;\cos\theta_t\right],\quad
    \theta_t = \tfrac{2\pi m_t}{1440}.
    \label{eq:meal-base-features}
\end{equation}
Here $G_t$ is the current glucose reading, $\Delta G_{\tau,t}$ are $\tau$-minute backward
slopes, and $(\sin\theta_t, \cos\theta_t)$ encode time of day continuously.

The retrospective setups build on this with increasingly rich leakage diagnostics:
\begin{equation}
    \underbrace{X_t^{\mathrm{base}}}_{A'}
    \;\longrightarrow\;
    \underbrace{\bigl[X_t^{\mathrm{base}},\,B_t\bigr]}_{B}
    \;\longrightarrow\;
    \underbrace{\bigl[X_t^{\mathrm{base}},\,\bar{p}_t^T\bigr]}_{C}.
    \label{eq:meal-retro-chain}
\end{equation}
The leakage-free chain replaces retrospective signals with causal ones:
\begin{equation}
    \underbrace{X_t^{\mathrm{base}}}_{A'}
    \;\longrightarrow\;
    \underbrace{\bigl[X_t^{\mathrm{base}},\,p_t^{T,\mathrm{causal}}\bigr]}_{C'}
    \;\longrightarrow\;
    \underbrace{\bigl[X_t^{\mathrm{base}},\,\widehat{p}_t^S\bigr]}_{D}
    \;\longrightarrow\;
    \underbrace{\bigl[X_t^{\mathrm{base}},\,\widehat{p}_t^S,\,\boldsymbol{\alpha}_t,\,\tau_t,\,c_t,\,\mathbf{q}_t\bigr]}_{H}.
    \label{eq:meal-causal-chain}
\end{equation}
The central empirical finding is:
\begin{equation}
    C \;\gg\; C' \;\approx\; D \;\approx\; H \;\approx\; A'.
    \label{eq:meal-central-comparison}
\end{equation}
The large performance advantage of C disappears when future-glucose access is removed; the
causal teacher, student, and online state carry little forecasting information beyond what
glucose level, slope, and clock already provide.

Table~\ref{tab:meal-setup-definitions} gives the formal definition of each setup's added
information, its mathematical representation, and its role in the evaluation.

\begin{table}[t]
\centering
\footnotesize
\renewcommand{\arraystretch}{1.35}
\caption{Formal definitions of the seven evaluation setups. ``Future info?'' marks setups
that use glucose or wearable data after time $t$ when computing the meal feature value.
These setups are diagnostic only and are not deployable.}
\label{tab:meal-setup-definitions}
\setlength{\tabcolsep}{4pt}
\begin{tabular}{@{}l p{3.8cm} p{4.0cm} p{3.2cm} c@{}}
\toprule
Setup & Information added to $X_t^{\mathrm{base}}$ & Meal representation $M_t^{(m)}$ & Role & Future? \\
\midrule
\textbf{q50}
  & None. Original SSMCGM-Stream median used directly.
  & $\widehat{y}_{t+h}^{q50}$ (no correction applied)
  & End-to-end baseline before any residual correction.
  & No \\[2pt]

\textbf{A$'$}
  & Current glucose $G_t$, three backward slopes $\Delta G_{\tau,t}$, clock $(\sin,\cos)\theta_t$.
  & $X_t^{\mathrm{base}}$ per Eq.~\eqref{eq:meal-base-features}
  & Fair harness baseline; isolates how much error slope and clock already explain.
  & No \\[2pt]

\textbf{B}
  & Thresholded binary predmeal flag from the bidirectional teacher: $B_t = \mathbb{1}\!\left[\bar{p}_t^T > \tau\right]$.
  & $\bigl[X_t^{\mathrm{base}},\,B_t\bigr]$
  & Leakage diagnostic for the old binary representation; retrospective.
  & Yes \\[2pt]

\textbf{C}
  & Corrected continuous teacher probability $\bar{p}_t^T = \frac{1}{|\mathcal{W}_t|}\sum_{w\in\mathcal{W}_t}p_{t,w}^T$.
  & $\bigl[X_t^{\mathrm{base}},\,\bar{p}_t^T\bigr]$
  & Retrospective oracle ceiling. Strong performance reflects future-glucose leakage, not deployable signal.
  & Yes \\[2pt]

\textbf{C$'$}
  & Causal teacher score computed from $G_{\le t}$ only: $p_t^{T,\mathrm{causal}} = f_\theta^T(G_{\le t})$.
  & $\bigl[X_t^{\mathrm{base}},\,p_t^{T,\mathrm{causal}}\bigr]$
  & Leakage-free teacher ceiling. Gap $C \gg C'$ quantifies how much of C's gain was unreachable future leakage.
  & No \\[2pt]

\textbf{D}
  & Causal student probability $\widehat{p}_t^S = f_\phi(X_{\le t}^{\mathrm{CGM}}, W_{\le t}, S_i)$, using past/current CGM, wearables, clock, and static context.
  & $\bigl[X_t^{\mathrm{base}},\,\widehat{p}_t^S\bigr]$
  & Basic deployable meal representation; tests whether the student adds signal beyond $A'$.
  & No \\[2pt]

\textbf{H}
  & Full online state: student probability plus structured phase decoder outputs $(\boldsymbol{\alpha}_t, \tau_t, c_t, \mathbf{q}_t)$, where $\boldsymbol{\alpha}_t$ is the online phase distribution (none/onset/rising/peak/recovery), $\tau_t$ is time since onset, $c_t$ is confidence, $\mathbf{q}_t$ is response-size probability.
  & $\bigl[X_t^{\mathrm{base}},\,\widehat{p}_t^S,\,\boldsymbol{\alpha}_t,\,\tau_t,\,c_t,\,\mathbf{q}_t\bigr]$
  & Richest deployable representation; adds temporal organization to D, but incremental gain is small when components correlate with slope already in $X_t^{\mathrm{base}}$.
  & No \\
\bottomrule
\end{tabular}
\end{table}

\subsection{Forecasting Evaluation and Selection Gate}
\label{subsec:meal-forecasting-evaluation}
Meal features were evaluated as residual corrections on top of the cached 60-minute SSMCGM-Stream \texttt{q50} forecast,
\begin{equation}
    \yhat^{(m)}_{i,t+h} = \yhat^{q50}_{i,t+h} + r_m(C_{i,t}, M^{(m)}_{i,t}, h),
    \label{eq:meal-residual-correction}
\end{equation}
fit on the validation cache and evaluated on the participant-disjoint test cache. Alignment passed on 1,660,377 test rows with maximum target residual 0.500 mg/dL and same-segment fraction 1.000.

Whole-horizon MAE60 is not the gate. A 60-minute average washes out the post-prandial excursion, and the negative controls below confirm almost no timing-specific signal exists on that endpoint. The gate is re-anchored on the endpoints where a meal acts, meal-window MAE, one hour peak error, time to peak, hyperglycemia AUPRC, and recall at precision 0.80 for the event \(\max(y \text{ over the next hour}) > 180\) mg/dL. A feature is selected only if, on the non-insulin cohort, it improves at least two of these endpoints over the baseline and beats shuffled, time-shifted, and block-shuffled controls with a participant-bootstrap interval that excludes zero. The decision is computed in code, not asserted in prose.

The baseline is A', a within-harness residual model on CGM slope, level, and clock only. The harness cost of A' relative to untouched \texttt{q50} is small on the non-insulin cohort, -0.005 mg/dL on MAE60 and -0.344 mg/dL on peak error, so meal comparisons are made against A' rather than against \texttt{q50} to separate the harness penalty from feature value. A' already captures most of the recoverable post-prandial signal, reaching meal-window MAE 13.745 and peak error 10.549 versus 13.856 and 10.893 for \texttt{q50}.

\begin{table}[t]
    \centering
    \caption{Reachable ceiling and leakage diagnostics on the non-insulin cohort. A' is the slope and clock baseline, C' is the leakage-free causal teacher ceiling, C and B are bidirectional-teacher leakage diagnostics. Lower is better for MAE, peak, and time to peak, higher is better for AUPRC and recall.}
    \label{tab:meal-main-results}
    \resizebox{\textwidth}{!}{%
\begin{tabular}{lrrrrrr}
\toprule
Setup & Meal MAE & Peak err. & TTP err. & Hyper AUPRC & Recall@.80 & MAE60 \\
\midrule
q50 & 13.856 & 10.893 & 19.4 & 0.876 & 0.776 & 14.357 \\
A$'$: slope+level+clock & 13.745 & 10.549 & 20.1 & 0.876 & 0.776 & 14.351 \\
C$'$: causal teacher (ceiling) & 13.724 & 10.542 & 20.0 & 0.876 & 0.778 & 14.333 \\
C: bidir teacher (leakage) & 11.286 & 8.150 & 17.5 & 0.926 & 0.862 & 11.789 \\
B: predmeal flag (leakage) & 13.226 & 9.876 & 19.7 & 0.888 & 0.792 & 13.737 \\
\bottomrule
\end{tabular}
}

\end{table}

The leakage-free causal teacher C' is the reachable ceiling, and it sits essentially at A', meal-window MAE 13.724 and peak error 10.542. The old bidirectional teacher C reaches meal-window MAE 11.286 and peak error 8.150, but that gap is leakage, because the only difference between C and C' is whether the teacher window is allowed to read future glucose. The previously reported gain from C, from 14.846 to 12.460 on MAE60, was therefore an oracle effect rather than reachable headroom. The deployable features cluster at A', the student probability D reaching meal-window MAE 13.711 and peak error 10.541, and the online full state H reaching 13.710 and 10.523.

\begin{table}[t]
    \centering
    \caption{Deployable causal representations against A' on the gate endpoints.}
    \label{tab:meal-causal-results}
    \resizebox{\textwidth}{!}{%
\begin{tabular}{lrrrrrr}
\toprule
Setup & Meal MAE & Peak err. & TTP err. & Hyper AUPRC & Recall@.80 & MAE60 \\
\midrule
A$'$: slope+level+clock & 13.745 & 10.549 & 20.1 & 0.876 & 0.776 & 14.351 \\
C$'$: causal teacher (ceiling) & 13.724 & 10.542 & 20.0 & 0.876 & 0.778 & 14.333 \\
D: student prob. & 13.711 & 10.541 & 19.7 & 0.877 & 0.777 & 14.266 \\
G: online state & 13.712 & 10.524 & 19.7 & 0.877 & 0.776 & 14.268 \\
H: online full state & 13.710 & 10.523 & 19.7 & 0.877 & 0.776 & 14.263 \\
\bottomrule
\end{tabular}
}

\end{table}

The orthogonalization test explains the clustering. After regressing each feature on slope, level, and clock, the leakage-free causal teacher C' has partial correlation 0.027 with the \texttt{q50} residual, the student probability D has 0.147, while the bidirectional teacher C has 0.519. The large orthogonal signal belongs only to the leaky features, so the deployable features carry little information that the CGM slope does not already provide.

\begin{table}[t]
    \centering
    \caption{Partial-information test against the CGM slope. Partial correlation with the q50 residual after removing slope, level, and clock, and the change in gate endpoints from adding the orthogonalized feature to the slope-only model.}
    \label{tab:meal-partial-information}
    \resizebox{\textwidth}{!}{%
\begin{tabular}{llrrr}
\toprule
Setup & Feature & Partial corr. & $\Delta$ Meal MAE & $\Delta$ Peak \\
\midrule
B: predmeal flag (leakage) & predmeal\_flag & 0.278 & 0.498 & 0.673 \\
C$'$: causal teacher (ceiling) & cgmacros\_teacher\_probability\_causal & 0.027 & 0.019 & 0.001 \\
C: bidir teacher (leakage) & cgmacros\_teacher\_probability & 0.519 & 2.302 & 2.385 \\
D: student prob. & student\_meal\_probability & 0.147 & 0.025 & 0.008 \\
G: online state & student\_meal\_probability & 0.147 & 0.025 & 0.008 \\
H: online full state & student\_meal\_probability & 0.147 & 0.025 & 0.008 \\
\bottomrule
\end{tabular}
}

\end{table}

\subsection{Leakage Audits and Negative Controls}
\label{subsec:meal-leakage-controls}
The provenance table records, for each feature, the maximum CGM input offset relative to the anchor. The student probability and the online states have offset zero, while the bidirectional teacher and the \texttt{predmeal\_flag} have a positive offset because their value at an anchor is computed from a window that extends past the anchor. A hard assertion fails the run if any teacher probability used as a selectable input has a positive offset, which is why the default teacher mode is causal.

\begin{table}[t]
    \centering
    \caption{Feature provenance and leakage audit. A positive input offset means the feature can read CGM later than its anchor.}
    \label{tab:meal-feature-provenance}
    \resizebox{\textwidth}{!}{%
\begin{tabular}{lrll}
\toprule
Feature & Max input offset (steps) & Uses future glucose & Leakage free \\
\midrule
student\_meal\_probability & 0 & False & True \\
online\_state\_features & 0 & False & True \\
predmeal\_flag & 71 & True & False \\
cgmacros\_teacher\_probability (legacy\_bidirectional) & 71 & True & False \\
cgmacros\_teacher\_probability\_causal (teacher\_mode=causal) & 0 & False & True \\
\bottomrule
\end{tabular}
}

\end{table}

The negative controls are the headline evidence. They compare the real feature timing against shuffled, time-shifted, and within-participant block-shuffled timing on the gate endpoints (Fig.~\ref{fig:meal-negative-controls}). For the deployable student probability D, real timing reaches peak error 10.541 against 10.601 shuffled, 10.541 shifted, and 10.500 block-shuffled, so the timing-specific gap is small. The leaky \texttt{predmeal\_flag} B shows a large real-versus-control gap, which is consistent with its bidirectional provenance.

\begin{table}[t]
    \centering
    \caption{Meal-state negative controls on the gate endpoints. Real timing versus shuffled, shifted, and block-shuffled meal features.}
    \label{tab:meal-negative-controls}
    \resizebox{\textwidth}{!}{%
\begin{tabular}{llrrr}
\toprule
Setup & Control & Meal MAE & Peak err. & Hyper AUPRC \\
\midrule
C$'$: causal teacher (ceiling) & real & 13.724 & 10.542 & 0.876 \\
C$'$: causal teacher (ceiling) & shuffle & 13.830 & 10.492 & 0.876 \\
C$'$: causal teacher (ceiling) & time\_shift & 13.859 & 10.716 & 0.874 \\
C$'$: causal teacher (ceiling) & block\_shuffle & 13.824 & 10.521 & 0.875 \\
D: student prob. & real & 13.711 & 10.541 & 0.877 \\
D: student prob. & shuffle & 13.792 & 10.601 & 0.876 \\
D: student prob. & time\_shift & 13.741 & 10.541 & 0.876 \\
D: student prob. & block\_shuffle & 13.813 & 10.500 & 0.876 \\
H: online full state & real & 13.710 & 10.523 & 0.877 \\
H: online full state & shuffle & 13.819 & 10.554 & 0.876 \\
H: online full state & time\_shift & 17.159 & 16.906 & 0.757 \\
H: online full state & block\_shuffle & 13.834 & 10.468 & 0.876 \\
B: predmeal flag (leakage) & real & 13.226 & 9.876 & 0.888 \\
B: predmeal flag (leakage) & shuffle & 14.116 & 14.174 & 0.866 \\
B: predmeal flag (leakage) & time\_shift & 13.714 & 10.689 & 0.868 \\
B: predmeal flag (leakage) & block\_shuffle & 14.206 & 12.196 & 0.871 \\
\bottomrule
\end{tabular}
}

\end{table}

\begin{figure}[t]
    \centering
    \includegraphics[width=\textwidth]{figures/meal_flags/negative_controls.pdf}
    \caption{Negative controls are the central evidence. The real-versus-control gap is the only honest measure of timing-specific meal signal, and it is small for the deployable features on the gate endpoints.}
    \label{fig:meal-negative-controls}
\end{figure}

\begin{table}[t]
    \centering
    \caption{Coded selection decision on the non-insulin cohort. A feature is selected only if it is leakage free and passes at least two gate endpoints against A' with controls beaten.}
    \label{tab:meal-selection}
    \begin{tabular}{lrll}
\toprule
Setup & Passing endpoints & Leakage free & Selected \\
\midrule
B: predmeal flag (leakage) & 5/2 & False & \textbf{False} \\
C$'$: causal teacher (ceiling) & 1/2 & True & \textbf{False} \\
C: bidir teacher (leakage) & 5/2 & False & \textbf{False} \\
D: student prob. & 1/2 & True & \textbf{False} \\
G: online state & 2/2 & True & \textbf{True} \\
H: online full state & 2/2 & True & \textbf{True} \\
\bottomrule
\end{tabular}

\end{table}

\subsection{Subgroup and Insulin Sensitivity Analyses}
\label{subsec:meal-subgroups}
Insulin users are quarantined from the gate and reported separately. On the non-insulin cohort the online full state H has peak error 10.523 versus 10.549 for A', a small difference. On insulin users the same comparison has much larger error and wider uncertainty, peak error 13.508 versus 13.456 for A', and MAE60 20.704 versus 20.805. Meal states in insulin users are therefore reported as low-confidence meal-like metabolic events, not literal meal detections.

\begin{table}[t]
    \centering
    \caption{Subgroup results split by insulin status for A', the online full state H, and the student probability D.}
    \label{tab:meal-subgroup-results}
    \resizebox{\textwidth}{!}{%
\begin{tabular}{llrrrr}
\toprule
Insulin & Setup & $n$ & Meal MAE & Peak err. & MAE60 \\
\midrule
med\_insulin=0 & q50 & 200 & 13.856 & 10.893 & 14.357 \\
med\_insulin=1 & q50 & 16 & 16.745 & 13.733 & 20.660 \\
med\_insulin=0 & A$'$: slope+level+clock & 200 & 13.745 & 10.549 & 14.351 \\
med\_insulin=1 & A$'$: slope+level+clock & 16 & 16.560 & 13.456 & 20.805 \\
med\_insulin=0 & D: student prob. & 200 & 13.711 & 10.541 & 14.266 \\
med\_insulin=1 & D: student prob. & 16 & 16.698 & 13.528 & 21.059 \\
med\_insulin=0 & H: online full state & 200 & 13.710 & 10.523 & 14.263 \\
med\_insulin=1 & H: online full state & 16 & 16.629 & 13.508 & 20.704 \\
\bottomrule
\end{tabular}
}

\end{table}

\begin{figure}[t]
    \centering
    \includegraphics[width=\textwidth]{figures/meal_flags/insulin_sensitivity.pdf}
    \caption{Insulin versus non-insulin split. The non-insulin story is clean, while uncertainty expands for insulin users.}
    \label{fig:meal-insulin-sensitivity}
\end{figure}

\subsection{Discussion and Limitations}
\label{subsec:meal-discussion-limitations}
The coded gate selected the following leakage free representations on the non insulin cohort: G on meal-window MAE and hyperglycemia AUPRC; H on meal-window MAE and hyperglycemia AUPRC. These clear the two endpoint bar against A' with shuffled, shifted, and block-shuffled controls beaten, so they may be described as useful on those endpoints. The effect sizes are small, on the order of a few hundredths of a mg/dL on meal-window MAE and a few thousandths on AUPRC, and the one hour peak error improvement does not survive all controls, so the result is a narrow, controlled gain rather than a large one. The reachable causal ceiling C' sits at the slope and clock baseline A', so the apparent gap to the bidirectional teacher is mostly unreachable leakage, not headroom. The partial-information test confirms that the deployable features carry little signal orthogonal to the CGM slope, which is expected because the weak labels are glucose-defined and the slope is already an input to \texttt{q50}. The practical reading is therefore cautious. The selected online states pass the gate only on meal-window MAE and hyperglycemia AUPRC, and only by small margins, while the larger one hour peak error improvement does not survive all three controls, so this is not yet evidence strong enough to justify full SSMCGM retraining. The shape-aware decoder adapter is reported as an ablation only and does not change this conclusion.

\begin{figure}[t]
    \centering
    \includegraphics[width=\textwidth]{figures/meal_flags/retrospective_causal_gap.pdf}
    \caption{Peak error and hyperglycemia AUPRC are the primary axes, with MAE60 demoted to the inset. The reachable causal ceiling C' sits at A', while the C leakage bar is unreachable headroom.}
    \label{fig:meal-retrospective-causal-gap}
\end{figure}

At the current checkpoint the meal-state branch is a validated weak-supervision and diagnostic pathway, not a selected production meal representation for SSMCGM retraining. The next step that could change the decision is a label that carries pre-rise information the CGM slope lacks, evaluated with the same coded gate.

\subsection{Pre-Rise Wearable Detector}
\label{subsec:prerise-wearable-detector}

To evaluate whether wearable signals can provide the missing lead time, we trained three causal feature sets on the binary \texttt{onset\_within\_k} target ($k = 6$ steps, 30 minutes) using a \texttt{HistGradientBoostingClassifier} with participant-disjoint validation. The feature sets are: \texttt{wearable\_only} (steps, heart rate, stress, sleep, clock; no raw CGM), \texttt{shape\_only} (a Kolle-style CGM shape score from a causal trailing window plus clock), and \texttt{wearable\_plus\_shape} (both). Hard leakage assertions prevent any future glucose or future wearable variable from entering the feature matrix. The shape score is permitted because it is computed from a causal trailing CGM window, not from raw level or slope. CGM-only baselines use a two-of-three ROC filter and a Harvey-style grid trigger applied to the same CGM series.

\begin{table}[t]
    \centering
    \caption{Pre-rise detection latency by setup on the non-insulin cohort ($n = 9{,}667$ onsets). Median detection latency $\Delta t$ relative to onset (negative = pre-rise); pre-rise fraction; sensitivity at threshold; false positives per day. Bootstrap 95\% CI from participant resampling.}
    \label{tab:prerise-detection-latency}
    \begin{tabular}{lrrrr}
        \toprule
        Setup & Median $\Delta t$ (min) & Pre-rise & Sensitivity & FP/day \\
        \midrule
        \texttt{wearable\_plus\_shape} & $0$   & 46.8\% & 0.81 & 11.5 \\
        \texttt{shape\_only}          & $20$  & 32.3\% & 0.70 &  7.8 \\
        \texttt{wearable\_only}       & $15$  & 35.7\% & 0.56 &  4.9 \\
        \texttt{D\_student\_prob}     & $10$  & 32.4\% & 0.12 &  0.4 \\
        \texttt{G\_online\_state}     & $20$  & 22.9\% & 0.63 &  2.1 \\
        \texttt{H\_online\_full\_state} & $20$ & 22.9\% & 0.63 & 2.1 \\
        CGM ROC 2-of-3 (baseline)    & $35$  & 18.3\% & 0.96 &  3.8 \\
        CGM grid-style (baseline)    & $40$  & 18.2\% & 0.99 &  3.6 \\
        \bottomrule
    \end{tabular}
\end{table}

Detection latency results are shown in Table~\ref{tab:prerise-detection-latency}. CGM-only baselines fire 35--40 minutes after onset on median, with pre-rise fractions below 19\%. The wearable feature sets fire earlier: \texttt{wearable\_only} achieves a 15-minute lead with 35.7\% pre-rise; \texttt{wearable\_plus\_shape} achieves a 0-minute median (half of detections are pre-rise) with 46.8\% pre-rise and sensitivity 0.81. These compare favorably to the Hochsmann et al.\ 37-minute CGM-only wall \citep{hochsmann2022detection}. The cost is a higher false-positive rate: \texttt{wearable\_plus\_shape} reaches 11.5 FP/day versus 3.6--3.8 for CGM baselines.

\paragraph{Forecasting gate.}
All setups fail the \texttt{selected\_and\_meaningful} gate. The MAE gain on meal windows is negative for all wearable feature sets (the detector fires before the glucose excursion, so the forecast window does not yet carry a recoverable residual), and none clear the 0.5 mg/dL MAE floor or the latency floor simultaneously. The CGM baselines pass the latency floor in the opposite direction (post-hoc, not a lead signal). This is an expected asymmetry, not a sign that the detector is wrong: a detector that fires before the CGM moves cannot improve a CGM-slope-based forecaster on that same window because the predictive information is not yet in the CGM. The correct metric for the wearable detector is detection latency, not meal-window MAE.

The conclusion mirrors the main pipeline: positive early-warning signal, no selected improvement over the coded forecasting gate. The wearable detector is reported as an early-warning diagnostic. The next evaluation asks whether it is ready to drive scenario counterfactuals through the SSMCGM decoder channel (Section~\ref{subsec:cf_phase0_results}).

% =========================================================
\subsection{Scenario-Aware Counterfactual Readiness}
\label{subsec:cf_scenario_retrain}
% =========================================================

\paragraph{Mask bug.}
An initial audit of the counterfactual path revealed that \texttt{\_apply\_proxy\_scenario} in the streaming evaluator unconditionally wrote \texttt{raw = 1.0} into the scenario channel regardless of the caller-supplied assertion, so every ``meal-off'' forward pass silently carried an active meal signal. The base model trained without scenario supervision had no mechanism to distinguish the two states, and the channel was functionally dead: validation factual-future gain $G_{\mathrm{ff}} \approx -0.0037$, indistinguishable from the no-signal baseline.

\paragraph{Scenario-aware retrain.}
Both arms were retrained from a shared five-epoch base checkpoint using a composite loss: factual-future pinball for correctly-asserted scenario windows, plus an event-window gain hinge with $\lambda_{\mathrm{gain}} = 0.1$ that rewards $G_{\mathrm{ff}} > 0$ on detected non-insulin event windows. The glucose-defined arm used the corrected predmeal flag (mean 0.1587, std 0.1346; raw 1.0 transforms to model-space 6.25). The wearable arm used the binary pre-rise score (mean 0, std 1, identity transform; raw 1.0 stays at model-space 1.0).

\paragraph{Channel revival.}
Both arms successfully revived $G_{\mathrm{ff}}$ from the dead-channel baseline. At epoch 8, the glucose-defined arm reached validation $G_{\mathrm{ff}} = {+0.64}$ and the wearable arm reached ${+0.42}$, both well above the ${-0.0037}$ baseline (Fig.~\ref{fig:f3_channel_revival}). A MAE60 forecast guard prevented checkpoint selection from overfitting: the wearable arm stayed below the guard limit throughout; the glucose arm showed forecast degradation after epoch 11, making epoch 8 the selected checkpoint for both arms. Positive $G_{\mathrm{ff}}$ is a necessary condition for meaningful counterfactual inference, but it does not alone establish that the model-implied effect is a valid causal estimate. The Block 3 audit below assesses that.

\begin{figure}[H]
    \centering
    \includegraphics[width=\textwidth]{figures/meal_flags/f3_channel_revival.pdf}
    \caption{F3: Scenario-channel revival across both arms. \textbf{Top:} Validation meal $G_{\mathrm{ff}}$ by epoch; both arms recover from the dead-channel baseline (dashed gray). Glucose arm (crimson) reaches $+0.64$, wearable arm (teal) reaches $+0.42$ at epoch 8. \textbf{Bottom:} Validation MAE60 forecast-only guard; wearable arm stays below the limit; glucose arm exceeds it after epoch 11, confirming epoch 8 as the selected checkpoint for both.}
    \label{fig:f3_channel_revival}
\end{figure}

% =========================================================
\subsection{Two-Arm Block 3 Counterfactual Evaluation}
\label{subsec:cf_block3_ab}
% =========================================================

Block 3 evaluates the 60-minute counterfactual effect ($\Delta\theta_{60}$) on the non-insulin cohort ($n = 200$ participants) at two operating points drawn from the pre-rise wearable detector: \texttt{wearable\_plus\_shape} (6592 anchors, 11.5 FP/day) and \texttt{wearable\_only} (4547 anchors, 4.9 FP/day). The natural-experiment reference $\Delta_{\mathrm{obs}}$ is $+22.2$ mg/dL at 60 minutes (bootstrap 95\% CI: 21.6 to 22.8), derived from matched non-insulin episodes observed to start and not start in the 30-minute window. A valid counterfactual should produce $\Delta\theta_{60}$ with the same sign as $\Delta_{\mathrm{obs}}$.

\paragraph{Glucose-defined arm.}
Both operating points yield strongly inverted counterfactual effects: $\Delta\theta_{60} = -35.3$ mg/dL (\texttt{wearable\_plus\_shape}, 95\% CI: $-36.2$ to $-34.4$) and $-35.1$ mg/dL (\texttt{wearable\_only}, 95\% CI: $-36.1$ to $-34.2$). The glucose-defined label selects anchors where CGM is already in a pre-excursion shape, which is exactly the set where the forecast-only path already anticipates a rise. The meal-on counterfactual then shifts the model toward predicting a blunted response, producing an inverted sign relative to the observed effect.

\paragraph{Wearable arm.}
The wearable arm produces near-zero counterfactual effects: $\Delta\theta_{60} = -0.20$ mg/dL (\texttt{wearable\_plus\_shape}, 95\% CI: $-0.212$ to $-0.197$) and $-0.21$ mg/dL (\texttt{wearable\_only}, 95\% CI: $-0.219$ to $-0.204$). Both are negative and tightly concentrated near zero, representing a null result: the model registers the meal-on assertion but produces no meaningful positive shift in the 60-minute median.

\paragraph{Dose-matched diagnostic.}
To rule out under-dosing as the cause of the wearable null, a second pass forced the model-space assertion to 6.25, matching the effective dose of the glucose-defined arm. At this dose, $\Delta\theta_{60} = -0.77$ mg/dL (\texttt{wearable\_plus\_shape}) and $-0.80$ mg/dL (\texttt{wearable\_only}). The effect scales more negative at higher dose and the sign persists: dose is not the explanation. Note that model-space 6.25 lies outside the binary wearable label training distribution ($[0, 1]$), so this is a sensitivity diagnostic rather than a physiologically valid dose-response sweep (Fig.~\ref{fig:f5_dose_sensitivity}).

\paragraph{Glycemic-state gradient.}
Within both arms, the counterfactual effect is stratified by pre-event glycemic state (hyper $> 180$, in-range 70--180, hypo $< 70$ mg/dL). Glucose arm at \texttt{wearable\_plus\_shape}: hyper $-39.4$ (n=255), in-range $-35.3$ (n=6256), hypo $-24.1$ (n=81). Wearable arm: hyper $-0.30$, in-range $-0.20$, hypo $-0.11$. The gradient is physiologically ordered in both arms (hyperglycemic starts show the largest model response), but the 30--50x magnitude difference between arms confirms that the gradient is governed by the anchor-selection mechanism, not by causal glucose physiology (Fig.~\ref{fig:f4_cf_ab_headline}).

\begin{figure}[H]
    \centering
    \includegraphics[width=\textwidth]{figures/meal_flags/f4_cf_ab_headline.pdf}
    \caption{F4: Counterfactual A/B headline. \textbf{Left:} $\Delta\theta_{60}$ at both operating points; glucose arm (crimson) is strongly inverted at approximately $-35$ mg/dL, wearable arm (teal) is near-zero null at approximately $-0.21$ mg/dL; both fail sign agreement with the observed $+22.2$ mg/dL (navy dashed). \textbf{Right:} Glycemic-state gradient; the ordering is preserved in both arms but the magnitude differs by 30 to 50x (y-axis scales differ).}
    \label{fig:f4_cf_ab_headline}
\end{figure}

\begin{figure}[H]
    \centering
    \includegraphics[width=0.65\textwidth]{figures/meal_flags/f5_dose_sensitivity.pdf}
    \caption{F5: Dose-sensitivity diagnostic. Wearable-arm $\Delta\theta_{60}$ at native model-space dose 1.0 and at dose 6.25 (matched to glucose-defined arm). Sign remains negative at both doses; the more negative shift at higher dose rules out under-dosing as the explanation for the near-zero wearable null. Dose 6.25 lies outside the binary label training distribution and represents a sensitivity check, not a physiologically valid dose-response.}
    \label{fig:f5_dose_sensitivity}
\end{figure}

% =========================================================
\section{Study 1: Meal Counterfactual Audit}
\label{sec:study1_cf_audit}
% =========================================================

This section collects the complete artifact-grounded evidence summary for the meal counterfactual audit. All values in figures and tables are read from CSV artifacts and cross-checked at generation time. No values are hard-coded in prose or scripts without a corresponding CSV row.

\subsection{Leakage and Redundancy}
\label{subsec:study1_leakage}

Figure~\ref{fig:f1_leakage_redundancy} and Table~\ref{tab:t1_ablation} report the seven-setup ablation. The central empirical result is Eq.~\eqref{eq:meal-central-comparison}: $C \gg C' \approx D \approx H \approx A'$. The leakage-free causal teacher C' sits at the slope-and-clock baseline A', while the bidirectional teacher C reaches substantially lower MAE and peak error, but that gap is future-glucose leakage and is fully unreachable at deployment. The partial correlation of the deployable student D with the q50 residual after removing slope, level, and clock is 0.101, compared to 0.519 for the leaky teacher C and 0.027 for the causal ceiling C'. The near-collapse of causal partial correlations confirms that the deployable features carry little information orthogonal to the CGM slope already available in the base model.

\begin{figure}[H]
    \centering
    \includegraphics[width=\textwidth]{figures/meal_flags/f1_leakage_redundancy.pdf}
    \caption{F1: Leakage and redundancy diagnosis. \textbf{Left:} Meal-window MAE for key setups; the causal ceiling C' equals the slope-and-clock baseline A'; the leaky oracle C gap (dashed crimson) is unreachable. \textbf{Center:} Partial correlation with q50 residual after removing slope, level, and clock; only the leaky teacher C (0.519) carries substantial orthogonal signal. \textbf{Right:} Negative controls for B (leaky) and C' (causal) on peak error; B shows a large real-versus-control gap consistent with bidirectional access to future glucose; C' does not.}
    \label{fig:f1_leakage_redundancy}
\end{figure}

\begin{table}[H]
\centering
\footnotesize
\renewcommand{\arraystretch}{1.3}
\caption{T1: Seven-setup ablation on the non-insulin cohort. MAE values in mg/dL. Partial correlation is with the q50 residual after partialing out CGM slope, level, and clock. ``Leaky'' marks setups that read future glucose or future wearables at the anchor. Dashes indicate no partial-correlation test was run for the baseline setups.}
\label{tab:t1_ablation}
\setlength{\tabcolsep}{5pt}
\begin{tabular}{@{}lllcccc@{}}
\toprule
Setup & Label & Feature added & Leaky & MAE60 & Meal-wdw MAE & Partial corr \\
\midrule
q50 uncorrected & Baseline & None (pretrained q50) & No & 14.36 & 13.86 & -- \\
A' & Slope+clock & CGM slope and clock & No & 14.35 & 13.74 & -- \\
B & Old predmeal & Binary predmeal flag (bidir) & Yes & 13.74 & 13.23 & 0.278 \\
C & Bidir teacher & Continuous bidir teacher prob & Yes & 11.79 & 11.29 & 0.519 \\
C' & Causal teacher & Causal teacher prob & No & 14.33 & 13.72 & 0.027 \\
D & Student prob & Causal student probability & No & 14.31 & 13.72 & 0.101 \\
G & Online state & Online decoded meal state & No & 14.30 & 13.71 & 0.101 \\
H & Online full & Online state + phase decoder & No & 14.28 & 13.71 & 0.101 \\
\bottomrule
\multicolumn{7}{@{}l}{\footnotesize Source: \texttt{block3\_meal\_feature\_endpoints.csv} (non\_insulin cohort),
  \texttt{block3\_partial\_information.csv}.}
\end{tabular}
\end{table}

\subsection{Pre-Rise Detection Performance}
\label{subsec:study1_prerise}

The pre-rise wearable detector evaluation is summarized in Fig.~\ref{fig:f2_prerise_detector}. The key metric is detection latency relative to CGM-defined onset: the CGM-only baselines fire 35--40 min post-onset, well above the Hochsmann et al.\ 37-min wall; the wearable feature sets fire at 0--15 min median, with \texttt{wearable\_plus\_shape} achieving a 0-min median (46.8\% of detections are pre-rise, sensitivity 0.81) at the cost of 11.5 FP/day. This confirms the detector breaks the CGM-only wall, motivating its use as the anchor-selection operating point for the Block 3 evaluation.

\begin{figure}[H]
    \centering
    \includegraphics[width=\textwidth]{figures/meal_flags/f2_prerise_detector.pdf}
    \caption{F2: Pre-rise wearable detector. \textbf{Left:} Median detection latency by setup; dotted red line marks the Hochsmann et al.\ 37-min CGM-only reference. Wearable sets (teal) fire at 0--15 min; CGM baselines (crimson diamonds) at 35--40 min. \textbf{Right:} Sensitivity vs.\ FP/day tradeoff; early detection of meal onset requires accepting a higher false-positive rate.}
    \label{fig:f2_prerise_detector}
\end{figure}

\subsection{Evidence Boundary and Conclusions}
\label{subsec:study1_boundary}

Table~\ref{tab:t2_block3} and Fig.~\ref{fig:f6_summary_schematic} summarize the Block 3 results. Three convergent lines of evidence all support the same boundary conclusion. First, the orthogonalization audit (Section~\ref{subsec:study1_leakage}) shows that deployable features carry no meal-specific signal orthogonal to the CGM slope, ruling out any headroom from better feature engineering. Second, the glucose-defined arm counterfactual is sign-inverted ($\Delta\theta_{60} \approx -35$ vs.\ observed $+22.2$ mg/dL), ruling out the glucose-defined label as a valid causal anchor. Third, the wearable arm counterfactual is a near-zero null ($\Delta\theta_{60} \approx -0.21$), and the dose-matched sweep to model-space 6.25 yields $-0.77$, ruling out under-dosing as an explanation. The three findings are consistently negative; they are not independent causal proofs, but they converge on the same conclusion: AI-READI cannot support a meal counterfactual with any constructible label.

\begin{table}[H]
\centering
\footnotesize
\renewcommand{\arraystretch}{1.3}
\caption{T2: Two-arm Block 3 counterfactual evaluation at 60 min horizon on the non-insulin cohort ($n = 200$ participants). $\Delta\theta_{60}$: model-implied meal-on minus forecast-only median; 95\% CI from participant bootstrap ($B = 1000$). $\Delta_{\mathrm{obs}} = +22.2$ mg/dL (95\% CI: 21.6 to 22.8) from matched natural experiment. Sign agreement column indicates whether $\Delta\theta$ and $\Delta_{\mathrm{obs}}$ share the same sign. Dose-matched rows use model-space 6.25 for the wearable arm (outside $[0,1]$ training distribution; diagnostic only).}
\label{tab:t2_block3}
\setlength{\tabcolsep}{4pt}
\begin{tabular}{@{}llrrrrrc@{}}
\toprule
Arm & Operating point & Dose & $\Delta\theta_{60}$ & 95\% CI & Anchors & Participants & Sign agree \\
\midrule
\multicolumn{8}{@{}l}{\textit{Glucose-defined arm (model-space 6.25 from preprocessor transform)}} \\
Glucose-def & \texttt{wearable\_plus\_shape} & 6.25 & $-35.3$ & $[-36.2,\ -34.4]$ & 6592 & 200 & No \\
Glucose-def & \texttt{wearable\_only} & 6.25 & $-35.1$ & $[-36.1,\ -34.2]$ & 4547 & 200 & No \\
\midrule
\multicolumn{8}{@{}l}{\textit{Wearable pre-rise arm (native dose, model-space 1.0)}} \\
Wearable & \texttt{wearable\_plus\_shape} & 1.0 & $-0.205$ & $[-0.212,\ -0.197]$ & 6592 & 200 & No \\
Wearable & \texttt{wearable\_only} & 1.0 & $-0.212$ & $[-0.219,\ -0.204]$ & 4547 & 200 & No \\
\midrule
\multicolumn{8}{@{}l}{\textit{Wearable arm (dose-matched to glucose-defined arm; sensitivity diagnostic only)}} \\
Wearable (dose-match) & \texttt{wearable\_plus\_shape} & 6.25 & $-0.775$ & $[-0.801,\ -0.747]$ & 6592 & 200 & No \\
Wearable (dose-match) & \texttt{wearable\_only} & 6.25 & $-0.798$ & $[-0.825,\ -0.771]$ & 4547 & 200 & No \\
\midrule
\multicolumn{8}{@{}l}{\textit{Glycemic-state gradient at \texttt{wearable\_plus\_shape} (glucose arm derived / wearable arm native)}} \\
Glucose-def & Hyper ($>180$ mg/dL) & 6.25 & $-39.4$ & $[-42.8,\ -36.4]$ & 255 & 41 & No \\
Glucose-def & In-range (70--180) & 6.25 & $-35.3$ & $[-36.2,\ -34.4]$ & 6256 & 199 & No \\
Glucose-def & Hypo ($<70$ mg/dL) & 6.25 & $-24.1$ & $[-29.9,\ -19.8]$ & 81 & 25 & No \\
Wearable & Hyper ($>180$ mg/dL) & 1.0 & $-0.303$ & -- & 255 & 41 & No \\
Wearable & In-range (70--180) & 1.0 & $-0.202$ & -- & 6256 & 199 & No \\
Wearable & Hypo ($<70$ mg/dL) & 1.0 & $-0.114$ & -- & 81 & 25 & No \\
\bottomrule
\multicolumn{8}{@{}l}{\footnotesize Sources: \texttt{block3\_60min\_operating\_point\_comparison.csv} (both arms),
  \texttt{block3\_60min\_glycemic\_state\_breakdown\_derived.csv} (glucose arm),}\\
\multicolumn{8}{@{}l}{\footnotesize \texttt{block3\_60min\_glycemic\_state\_breakdown.csv} (wearable arm),
  \texttt{natural\_experiment\_effects.csv}.}
\end{tabular}
\end{table}

\begin{figure}[H]
    \centering
    \includegraphics[width=\textwidth]{figures/meal_flags/f6_summary_schematic.pdf}
    \caption{F6: Three convergent evidence streams supporting the boundary conclusion. Each stream rules out one alternative explanation (leakage headroom, encoding artefact, dose artefact), and all three converge on the same verdict. The streams are independently consistent; they are not independent causal proofs. The boundary conclusion follows from their convergence, not from any single test.}
    \label{fig:f6_summary_schematic}
\end{figure}

\paragraph{Boundary conclusion.}
AI-READI cannot support a meal counterfactual in any constructible label given the current dataset. The observed glucose rise is $\Delta_{\mathrm{obs}} = +22.2$ mg/dL at 60 min. Both constructible labels fail sign agreement at every dose and every operating point. The failure is not attributable to dead channels (G\textsubscript{ff} was successfully revived), dose magnitude (dose-matched sweep confirms), or encoding artefacts (lead preservation was confirmed in Step 2). The required next step is exogenous meal annotation from a dataset with prospective logging (for example, T1DEXI or CGMacros), which would enable a label that carries pre-rise causal information not derivable from the CGM slope alone.

\subsection{Study Framing}
\label{subsec:study_framing}

This report presents one of three planned studies.

\textbf{Study 1: Meal Counterfactual Audit} (this section). Concludes that AI-READI does not support a meal counterfactual with any constructible label. The boundary is established by three convergent checks: orthogonalization, glucose-defined arm inversion, and wearable arm dose-scaled null. The recommended next step is evaluation on a dataset with exogenous meal annotations.

\textbf{Study 2: Insulin and Medication Counterfactuals.} Asks whether the scenario-aware channel supports a valid insulin-dose or medication-adherence counterfactual on AI-READI, where external covariates provide labeling that does not depend on the CGM slope. The insulin quarantine and the low-confidence medication flags established in this report define the scope of that evaluation.

\textbf{Study 3: Downstream Clinical Decision Support.} Asks whether the SSM-CGM streaming distribution, with or without a valid counterfactual channel, supports reliable hyperglycemia or hypoglycemia risk alerts at the operating points identified by the pre-rise detector. This study uses the hyperglycemia AUPRC and recall endpoints already reported in Table~\ref{tab:t1_ablation} as its primary metrics.
